\paragraph{The exact conjugate determinant comparison.}
The passage from local flat norms to the degrees of the prescribed
eigenlattices has a short direct proof. Let $K$ be a number field and
$\bar V$ an integral Hermitian lattice with an action of $\mathcal O_E$.
For two conjugate embeddings into $K$, let $W,W'$ be the integral
eigenlattices, of ranks $a,b$, with their inherited metrics. Put
$\mathcal D=\det W\otimes(\det W')^{-1}$.

\begin{lemma}\label{geo-flat-determinant-bridge}
After one common finite extension if necessary, suppose a scalar
$\varphi(x,y)\in K^\times$ attached to rational eigenframes commutes
with field embeddings and field extension, and satisfies
\[
 \varphi(Ax,By)=\varphi(x,y)\frac{\det B}{\det A}.
\]
At complex embeddings suppose its specified logarithmic norm is
\[
 \log|\varphi(x,y)|+
 \tfrac12\log\det\operatorname{Gram}(x)-
 \tfrac12\log\det\operatorname{Gram}(y).
\]
At every finite place suppose it is $\log|\varphi(x_0,y_0)|$
on integral local frames of these same eigenlattices. Then its normalized
adelic degree is
$\widehat{\deg}_{\rm n}W-\widehat{\deg}_{\rm n}W'$.
The statement remains valid after a common finite field extension,
including at ramified places.
\end{lemma}
\begin{proof}
The section
$s=\varphi(x,y)(\bigwedge x)\otimes(\bigwedge y)^{-1}$ of
$\mathcal D_K$ is independent of the frames. Its complex norm is the
displayed Gram expression. At a finite place, writing $x=Ax_0$ and
$y=By_0$ gives
\[
 \log\|s\|=\log|\varphi(x,y)|+\log|\det A|-\log|\det B|
                         =\log|\varphi(x_0,y_0)|.
\]
Thus the supplied norms are exactly those of $\mathcal D$. Sum their
negatives over all embeddings and finite places, and divide by
$[K:\mathbb Q]$. The product formula removes any change of rational
section. For integral base change, choose an integral basis $e_i$ of
$\mathcal O_E$: an eigenlattice is the kernel of
$v\mapsto((e_i-\tau(e_i))v)_i$. Flat extension preserves this kernel,
so no denominator or new index occurs. Normalized degree is unchanged
because the embeddings occur with their extension multiplicities.
For rank zero the determinant is the trivial line.
\end{proof}

The geometric input is a flat family with these local norms and their
compatibility with the prescribed complex multiplication points.
Proposition~\ref{geo-cm-flat-degree} below computes its value from one
elliptic curve under precisely that compatibility. The finite-place
normalization must be retained. Write
\[
 P_E=\gamma+\log(2\pi)-\log\operatorname{rd}(E)
                      -\frac{L'(1,\chi_{-d})}{L(1,\chi_{-d})}.
\]
Suppose the corrected flat degree is
$(a-b)P_E/2+b(\log d-\log\operatorname{rd}(E))$ and the actual
integral-frame logarithmic norms add
$b(\log|\mathfrak D_E|-\log|d|)$ at finite places. Their normalized
sum is $b(\log d-\log\operatorname{rd}(E))$, by the ideal norm
of the different and the product formula. Degree is the negative
sum of logarithmic norms, so this term cancels exactly. The lemma gives
\begin{equation}\label{geo-exact-conjugate-degree}
 \widehat{\deg}_{\rm n}W-\widehat{\deg}_{\rm n}W'
                         =\frac{a-b}{2}P_E.
\end{equation}
These are hypotheses about the same semistable model and its prescribed
metric. They do not follow from a scalar estimate for $P_E$.

\begin{corollary}\label{geo-signed-signature-budget}
In the notation of \eqref{geo-selected-index-identity}, assume the
flat-family and period inputs above for each conjugate pair with the
same $P_E$. Assume also the tangent-slope estimate
$T\le12R\log(2+n+(-d)_+)$ for these same lattices, with $d$ their
total degree. Then
\[
 b=\frac{2R-n}{2}P_E,\qquad
 T\le12R\log\left(2+n+
       \left(\frac{2R-n}{2}P_E-\delta-2\nu\right)_+\right).
\]
For constant signature $(p,q)$, toric rank $r$, and a selected half
of size $m\in\{p,q\}$, one has
$2R-n=f(2m-p-q)$, independently of $r$. Equal signatures remove
the period cost. Selecting the smaller half also removes it when
$P_E\ge0$.
\end{corollary}
\begin{proof}
Sum \eqref{geo-exact-conjugate-degree}; the complementary ranks sum
to $n-R$. Apply the selected-degree argument with the tangent-slope
constant $12$. Finally use $R=f(m-r)$ and $n=f(p+q-2r)$.
\end{proof}
